% ============================================================
%  宽松风险平价答辩 PPT
%  编译：xelatex rrp_defense.tex  →  rrp_defense.pdf
%  位置：report/defense_ppt/
%  图片：../../results/figures/（相对于本文件）
% ============================================================
\documentclass[aspectratio=169,11pt]{beamer}
\usepackage{beamerthemeSWUFE}
\usepackage{booktabs}
\usepackage{amsmath}
\usepackage{array}
\usepackage{colortbl}

% 图片搜索路径：results/figures 和本地 logo
\graphicspath{{../../results/figures/}{assets/logos/}}

% ── 封面信息 ─────────────────────────────────────────────────────────────────
\title{宽松风险平价在全球 ETF\\资产配置中的改进与实证研究}
\subtitle{本科毕业论文答辩}
\author{许哲圣}
\studentid{42353012}
\school{特拉华数据科学学院}
\major{金融数学（量化分析方向）}
\supervisor{王磊}
\date{2026年8月}

% ── 数值宏：与论文共用同一流水线生成文件 ───────────────────────────────────────
\newcommand{\evalStartDate}{2018-01-02}
\newcommand{\evalEndDate}{2026-09-11}
\newcommand{\etfCount}{30}
\newcommand{\txCostBps}{3}
\newcommand{\annualizationDays}{252}
\newcommand{\lookbackDays}{252}
\newcommand{\primaryRebalances}{445}
\newcommand{\primaryObservations}{2111}
\newcommand{\primaryCashMean}{18.02\%}
\newcommand{\primaryCashMax}{26.43\%}
\newcommand{\primaryDefensiveMean}{66.29\%}
\newcommand{\primaryAssetsUsed}{30}
\newcommand{\primarySelectionInformativeYears}{0}
\newcommand{\primaryParameterYears}{10}
\newcommand{\primaryReturnTargetDef}{$R_t=\mu_t^\mathsf{T}q_t$，即当期可行风险预算参考的预测收益}
\newcommand{\globalGrossReturn}{6.15\%}
\newcommand{\globalNetReturn}{6.09\%}
\newcommand{\globalVolatility}{3.59\%}
\newcommand{\globalMaxDD}{-6.80\%}
\newcommand{\globalMonthlyTurnover}{17.38\%}
\newcommand{\globalCostDrag}{0.07\%}
\newcommand{\globalCVaR}{0.53\%}
\newcommand{\globalSharpe}{1.664}
\newcommand{\globalSortino}{2.377}
\newcommand{\globalCalmar}{0.895}

\newcommand{\primaryCashMean}{70.23\%}
\newcommand{\primaryCashMax}{86.61\%}
\newcommand{\primaryChinaBondSharpe}{0.529}
\newcommand{\primaryRebalances}{444}
\newcommand{\primaryObservations}{2102}


% ── 辅助命令：图片帧（单图 + 可选注释）────────────────────────────────────────
% \figframe{帧标题}{文件名}{宽度比}{说明文字（空=不显示）}
\newcommand{\figframe}[4]{%
  \begin{frame}{#1}
    \begin{center}
      \includegraphics[width=#3\textwidth,height=0.62\textheight,keepaspectratio]{#2}
    \end{center}
    \begin{center}\scriptsize #4\end{center}
  \end{frame}%
}


\begin{document}
\begin{frame}[plain]\titlepage\end{frame}
\section{研究问题}
\begin{frame}{风险预算与实施约束}
\begin{itemize}
\item 30 只中国市场可交易 ETF，覆盖 8 类全球资产风险
\item 先求凸风险预算参考，再用凸跟踪与换手目标生成权重
\item 研究集中度约束是否实际改变配置和历史绩效
\item 主模型为 Improved Convex Adaptive Global RRP，其余模型为对照
\end{itemize}
\end{frame}
\begin{frame}{数据与评价口径}
\begin{itemize}
\item 评价区间：\evalStartDate{} 至 \evalEndDate{}
\item \primaryObservations{} 个日收益观察，243 日年化，无风险利率为零
\item 主模型每周最后一个实际交易日调仓，共 \primaryRebalances{} 次
\item 单边成本 \txCostBps{} bps，为固定研究约定
\item 保留极端收益，上市前不回填，至少 60 个历史观察才可投资
\end{itemize}
\end{frame}
\section{主模型方法}
\begin{frame}{凸优化主模型}
风险预算参考 $b_t$ 来自凸对数障碍问题。实施阶段求解
\[
\min_w\;0.25\|w-b_t\|_2^2+(0.02+0.0003)\|w-w_t^-\|_1
\]
\[
\mathbf{1}^{T}w=1,\quad w\geq0,\quad\|w-w_t^-\|_1\leq0.80.
\]
\begin{itemize}
\item 取消现金组与单资产集中度上限，不增加杠杆
\item 保留债券、防御、商品与黄金、权益组边界
\item 当前收益奖励、方差惩罚和 CVaR 惩罚均为零
\end{itemize}
\end{frame}
\begin{frame}{约束与参数来源}
\begin{tabular}{lr}
\toprule
参数 & 主模型设置\\\midrule
历史窗口 / EWMA 半衰期 & 252 / 60 个观察\\
债券组上限 & 70\%\\
防御组上限 & 25\%\\
商品与黄金组上限 & 40\%\\
权益组上限 & 70\%\\
单期换手上限 & 80\%\\\bottomrule
\end{tabular}
\medskip

这些参数继承自研究配置，未宣称为实测最优值。现金无额外上限，但受多头单位预算约束。
\end{frame}
\begin{frame}{时序与选择边界}
\begin{itemize}
\item 每期优化只使用严格早于调仓日的数据
\item 沿用当日目标权重乘当日收益并扣成本的日频记账约定
\item 冻结既有候选日历，基础候选为 \texttt{candidate\_03}
\item 周频与取消上限规格是在历史实验后选定的
\item 因此属于探索性历史证据，不是独立样本外模型选择验证
\end{itemize}
\end{frame}
\section{实证结果}
\begin{frame}{主模型绩效}
\begin{tabular}{lr}
\toprule
指标 & 结果\\\midrule
净年化收益 & \improvedNetReturn\\
年化波动 & \improvedVolatility\\
夏普（无风险利率为零） & \improvedSharpe\\
Sortino & \improvedSortino\\
最大回撤 & \improvedMaxDD\\
平均月度换手 & \improvedMonthlyTurnover\\\bottomrule
\end{tabular}
\medskip

月均换手按自然月汇总，实际调仓频率为周频。
\end{frame}
\begin{frame}{核心模型对照}
\scriptsize
\resizebox{\textwidth}{!}{\begin{tabular}{lrrrr}
\toprule
模型 & 净年化收益 & 年化波动 & 夏普 & 最大回撤\\\midrule
Improved Convex Adaptive Global RRP & \improvedNetReturn & \improvedVolatility & \improvedSharpe & \improvedMaxDD\\
Global RRP & \globalNetReturn & \globalVolatility & \globalSharpe & \globalMaxDD\\
Convex Adaptive Global RRP & \convexNetReturn & \convexVolatility & \convexSharpe & \convexMaxDD\\
HRP & \hrpNetReturn & \hrpVolatility & \hrpSharpe & \hrpMaxDD\\
HERC & \hercNetReturn & \hercVolatility & \hercSharpe & \hercMaxDD\\
Equal Weight & \equalWeightNetReturn & \equalWeightVolatility & \equalWeightSharpe & \equalWeightMaxDD\\
60/40 & \sixtyFortyNetReturn & \sixtyFortyVolatility & \sixtyFortySharpe & \sixtyFortyMaxDD\\\bottomrule
\end{tabular}}
\medskip

主模型为周频，其余核心对照为月频。主模型身份不表示每项指标排名第一。
\end{frame}
\figframe{累计净值}{convex_adaptive_nav_comparison.png}{0.88}{相同日期与成本口径，结合风险尺度理解绝对收益差异}
\figframe{历史回撤}{convex_adaptive_drawdown_comparison.png}{0.88}{主模型最大回撤为 \improvedMaxDD，未启用 CVaR 惩罚}
\figframe{换手与实施成本}{convex_adaptive_turnover_comparison.png}{0.88}{换手影响成本。固定 3 bps 未计入实测价差、容量与市场冲击}
\figframe{现金集中度}{primary_weights.png}{0.92}{平均现金类 ETF 权重 \primaryCashMean，最高 \primaryCashMax}
\begin{frame}{无风险利率与夏普解释}
同一组合、同一权重、同一净收益路径：
\begin{center}
\begin{tabular}{lr}
\toprule
无风险利率口径 & 夏普\\\midrule
固定为零 & \improvedSharpe\\
滞后一月的中债 1 年期国债收益率 & \primaryChinaBondSharpe\\\bottomrule
\end{tabular}
\end{center}
较高现金权重与极低波动会显著影响零利率比率。口径变化不创造额外收益，近现金对照也可能得到更高夏普。
\end{frame}
\section{诊断与局限}
\begin{frame}{约束消融}
\scriptsize
\resizebox{\textwidth}{!}{\begin{tabular}{lrrrr}
\toprule
生效年 & $\lambda_{v,t}$ & $\lambda_{r,t}$ & 验证夏普 & 标准误集合\\\midrule
2017 & 6.858e-12 & 16.54 & 1.467 & 9/9\\
2018 & 0.1402 & 8.826 & 1.630 & 9/9\\
2019 & 0.08214 & 24.72 & 0.026 & 9/9\\
2020 & 0.1012 & 9.385 & 3.433 & 9/9\\
2021 & 0.1181 & 2.936 & 1.629 & 9/9\\
2022 & 0.08515 & 2.01 & 2.051 & 9/9\\
2023 & 0.1336 & 8.52 & 0.280 & 9/9\\
2024 & 0.2022 & 11.52 & 2.796 & 9/9\\
2025 & 0.1931 & 15.49 & 2.547 & 9/9\\
2026 & 0.1574 & 3.793 & 2.207 & 9/9\\
\bottomrule
\end{tabular}}
\medskip

仅取消现金组上限仍保留单资产上限。相对 CVaR 保留原最优基准的可行性，微小数值差异不构成独立改善证据。
\end{frame}
\begin{frame}{分自然年结果}
\scriptsize
\centering
\begin{tabular}{lrrrr}
\toprule
年份 & 净年化收益 & 年化波动 & 夏普 & 最大回撤\\\midrule
2018 & 0.06\% & 3.79\% & 0.034 & -5.04\%\\
2019 & 14.29\% & 8.11\% & 1.696 & -4.30\%\\
2020 & 22.20\% & 10.86\% & 1.908 & -7.16\%\\
2021 & 9.24\% & 6.12\% & 1.481 & -3.37\%\\
2022 & -1.58\% & 3.95\% & -0.384 & -3.19\%\\
2023 & 3.93\% & 3.44\% & 1.144 & -1.76\%\\
2024 & 12.51\% & 10.86\% & 1.144 & -6.95\%\\
2025 & 16.30\% & 8.67\% & 1.793 & -3.34\%\\
2026 & 22.82\% & 8.35\% & 2.519 & -4.69\%\\
\bottomrule
\end{tabular}
\medskip

2026 年仅截至 \evalEndDate。年度年化指标不是完整年度已实现收益。
\end{frame}
\figframe{主模型参数下的频率对照}{rebalance_frequency_sensitivity.png}{0.86}{仅改变周、双周、月、季度日历。周频为指定配置，排名属于描述性结果}
\begin{frame}{期末持仓快照}
\centering
% Auto-generated by scripts/export_next_month_holdings.py
\scriptsize\renewcommand{\arraystretch}{1.12}
\begin{tabular}{llr}
\toprule
ETF & 代码 & 权重 \\
\midrule
日利ETF & 511880.SH & 79.94\% \\
信用债ETF & 511030.SH & 5.44\% \\
5年国债ETF & 511010.SH & 3.72\% \\
10年国债ETF & 511260.SH & 2.56\% \\
标普500ETF & 513500.SH & 0.75\% \\
纳指ETF & 159941.SZ & 0.75\% \\
红利ETF & 510880.SH & 0.71\% \\
原油ETF & 162411.SZ & 0.67\% \\
恒生ETF & 159920.SZ & 0.66\% \\
欧洲ETF & 513030.SH & 0.65\% \\
\bottomrule
\end{tabular}

\medskip

\footnotesize 展示期末最大持仓，不是下月整月固定指令。周频模型会在下一再平衡点重新求解。
\end{frame}
\begin{frame}{验证与研究边界}
\begin{itemize}
\item 核查共同日期、权重归一、漂移权重、成本、约束及信息截止日
\item 主模型路径复现获批研究结果，保留求解与约束审计
\item 高现金集中度限制了风险分散的解释
\item 事后规格选择尚需新增样本验证，历史 PBO 不自动适用于当前规格
\item 未建模负债、资本占用、容量与实际成交，不能直接推出机构适用性
\end{itemize}
\end{frame}
\section{结论}
\begin{frame}{研究结论}
\begin{itemize}
\item 周频凸优化主模型在零利率口径下达到 \improvedSharpe{} 夏普
\item 净收益为 \improvedNetReturn，平均现金类 ETF 权重为 \primaryCashMean
\item 约束消融区分结构变化与统计口径变化
\item 其余模型提供对照，未来验证应冻结规格并使用新增数据
\end{itemize}
\medskip

历史结果不保证未来收益或夏普达到同一水平。
\end{frame}
\begin{frame}[plain]\centering\Huge 谢谢各位老师\end{frame}
\end{document}
